\section{Lines in Three-Dimensional Space}

A point in three-dimensional Cartesian space has coordinates
\[
P=(x,y,z)\in\mathbb{R}^3.
\]
A line is determined by one point and one nonzero direction.

\begin{definitionbox}
Let
\[
P_0=(x_0,y_0,z_0)
\]
and let
\[
\mathbf{v}=(a,b,c)\neq(0,0,0).
\]
The line through $P_0$ with direction vector $\mathbf{v}$ is
\[
L=\{P_0+t\mathbf{v}:t\in\mathbb{R}\}.
\]
Equivalently,
\[
(x,y,z)=(x_0,y_0,z_0)+t(a,b,c).
\]
\end{definitionbox}

In coordinates this becomes the system
\[
\begin{cases}
x=x_0+at,\\
y=y_0+bt,\\
z=z_0+ct,
\end{cases}
\qquad t\in\mathbb{R}.
\]
As $t$ varies, the point moves along the line. Positive and negative values of
$t$ describe the two opposite directions along the same geometric line.

\begin{center}
\begin{tikzpicture}[scale=0.9]
    \draw[axis] (0,0) -- (5.2,0) node[right] {$x$};
    \draw[axis] (0,0) -- (-1.7,-1.5) node[below left] {$y$};
    \draw[axis] (0,0) -- (0,4.3) node[above] {$z$};

    \coordinate (P) at (1.2,1.0);
    \coordinate (Q) at (4.2,3.0);
    \draw[subset] (0.5,0.53) -- (4.8,3.4) node[above right] {$L$};
    \fill[geometricSubsetColor] (P) circle (2.2pt) node[below left] {$P_0$};
    \draw[->,thick,geometricSubsetColor] (P) -- (Q) node[midway,above left] {$\mathbf v$};
\end{tikzpicture}
\end{center}

\subsection{The Line Through Two Points in Space}

Let
\[
P=(x_1,y_1,z_1),
\qquad
Q=(x_2,y_2,z_2)
\]
be distinct points. Their difference
\[
Q-P=(x_2-x_1,\,y_2-y_1,\,z_2-z_1)
\]
is a nonzero direction vector of the line through them.

\begin{theorembox}
The unique line through two distinct points $P,Q\in\mathbb{R}^3$ is
\[
L=\{P+t(Q-P):t\in\mathbb{R}\}.
\]
In coordinates,
\[
\begin{cases}
x=x_1+t(x_2-x_1),\\
y=y_1+t(y_2-y_1),\\
z=z_1+t(z_2-z_1).
\end{cases}
\]
\end{theorembox}

\begin{examplebox}
For
\[
P=(1,0,2),
\qquad
Q=(3,4,1),
\]
a direction vector is
\[
Q-P=(2,4,-1).
\]
Therefore the line through $P$ and $Q$ is
\[
(x,y,z)=(1,0,2)+t(2,4,-1),
\]
or
\[
\begin{cases}
x=1+2t,\\
y=4t,\\
z=2-t.
\end{cases}
\]
\end{examplebox}

\subsection{Symmetric Form}

If $a$, $b$, and $c$ are all nonzero, the parameter can be eliminated from
\[
x=x_0+at,
\qquad
y=y_0+bt,
\qquad z=z_0+ct.
\]
This gives
\[
\frac{x-x_0}{a}
=
\frac{y-y_0}{b}
=
\frac{z-z_0}{c}.
\]
If one component of the direction vector is zero, the corresponding coordinate
is constant. For example, if $c=0$, then $z=z_0$, while the remaining two
coordinates may still be written in symmetric form.

\subsection{Parallel and Intersecting Lines in Space}

Consider
\[
L_1=P_1+t\mathbf v_1,
\qquad
L_2=P_2+s\mathbf v_2.
\]
The lines are parallel when their direction vectors are proportional:
\[
\mathbf v_2=\lambda\mathbf v_1
\]
for some nonzero $\lambda$.

To test whether two nonparallel lines intersect, one solves
\[
P_1+t\mathbf v_1=P_2+s\mathbf v_2
\]
for the two parameters $t$ and $s$. This is a system of three scalar equations.
It may have one solution or no solution.

\begin{definitionbox}
Two lines in $\mathbb{R}^3$ are \emph{skew} when they are neither parallel nor
intersecting.
\end{definitionbox}

Skew lines have no analogue in the plane. In two dimensions, two distinct
nonparallel lines must intersect; in three dimensions, one line may pass above
or beside another without meeting it.

\begin{summarybox}
A line in $\mathbb{R}^3$ is described by
\[
L=\{P_0+t\mathbf v:t\in\mathbb{R}\},
\qquad \mathbf v\neq\mathbf 0.
\]
Two distinct points determine the direction $Q-P$. Parallel lines have
proportional direction vectors, while nonparallel lines may either intersect or
be skew.
\end{summarybox}